% Source-based authoring: Girard's original trip criterion and a distinct modern
% binary unit-free switching proof.  Original DR1989 source completion is open.
\chapter{Multiplicative Proof Structures, Criteria, and Cut}
\label{sec:p07-multiplicative-proof-nets}

Proof structures record formula occurrences and their links, not just a
sequent's list of formulas. Correctness is a global condition on that
incidence data. We first keep Girard's original long-trip presentation
separate from a binary switching presentation, and then give a complete
sequentialization and cut argument for the latter. The modern splitting
argument is reconstructed from
\cite[pp.~1--5]{DiGuardiaLaurentTortoraVaux2023Sequentialization};
it is not attributed as the proof printed in
\cite{DanosRegnier1989}. Girard's actual original treatment is
\cite[Section~2, pp.~28--47]{Girard1987AuthorHosted}.

\section{Occurrences, ports, and the selected fragment}
\label{sec:p07-mll-objects}

\begin{hbdefinition}{def:p07-mll-formulas}{Unit-free multiplicative formulas}
Fix paired atomic symbols \(a,a^\perp\). Formulas are generated by those
symbols and the binary constructors \(\otimes,\bindnasrepma\). Linear
negation is the involution determined by
\[
 (a^\perp)^\perp=a,\qquad
 (A\otimes B)^\perp=A^\perp\bindnasrepma B^\perp,\qquad
 (A\bindnasrepma B)^\perp=A^\perp\otimes B^\perp.
\]
Occurrences are separately identified objects carrying these formula labels.
Equal labels do not identify occurrences. An axiom may join arbitrary
dual formulas \(A,A^\perp\); it is not silently restricted to atomic
axioms. Its two outgoing occurrences remain opaque at that link.
This is the unit-free fragment of the one-sided language in
\cref{def:p08-cll-language}. The units \(\mathbf1,\bot\), additives and
exponentials of \cref{sheet:cll-one-sided} are not constructors of this
selected net system.
\end{hbdefinition}

\begin{hbdefinition}{def:p07-mll-partial-graph}{The incidence object}
A proof structure consists of finite sets \(V,E\), partial endpoint maps
\(s,t:E\rightharpoonup V\), labels on arcs, and types
\(\mathsf{ax},\mathsf{ten},\mathsf{par}\) on vertices.
Arcs are formula occurrences. At an axiom vertex there are no incoming
arcs and two outgoing arcs labeled \(A,A^\perp\). A tensor or par vertex
has two ordered incoming arcs labeled \(A,B\) and one outgoing arc
labeled \(A\otimes B\) or \(A\bindnasrepma B\), respectively.
Each defined endpoint is a single endpoint, so an occurrence is not
consumed by two links. An arc with undefined source is a marked
hypothesis; an arc with undefined target is a conclusion.
An arc with both endpoints undefined is allowed in the auxiliary open
system. A closed structure has no hypothesis arcs.

For undirected connectivity, fill each undefined endpoint with its own
boundary point. In particular, an isolated open arc is one component,
not an empty graph. A path is a walk through incident arcs with a
specified direction on each traversal; a simple cycle is nonempty,
with no repeated vertex except its endpoint and no repeated arc.
Parallel arcs are distinct, so a two-arc cycle counts as a cycle.
Formula complexity increases from premises to conclusion at a tensor
or par link. Thus the directed link incidence has no directed cycle.
The object is a graph of occurrences, not a drawing with geometrical
crossings interpreted as vertices.
\end{hbdefinition}

\begin{hbdefinition}{def:p07-mll-switching}{Switching, acyclicity, and connectedness}
A switching chooses one premise incidence at every par vertex. Sever
the other incidence and give its premise arc a fresh boundary endpoint;
do not erase that formula occurrence. Tensor incidences and axiom arcs
are retained. Write \(G_\sigma(N)\) for the resulting undirected graph.

Call \(N\) \emph{switching-acyclic} if every \(G_\sigma(N)\) is a forest.
Call it \emph{switching-correct and connected} if every such graph is a
tree. The empty structure is acyclic, but is not connected-correct.
Equivalently, a switching cycle is a simple cycle that never uses both
premises of one par vertex, and acyclicity excludes every switching cycle.
Connectedness is an additional condition, not a consequence of acyclicity.
For an acyclic structure write \(d(N)\) for its number of switching
components. This is independent of the choice: every switching has the
same number of vertices and arcs, and a finite forest has component
count equal to its number of vertices minus its number of edges.
The acyclic auxiliary calculus permits the two Mix rules; the connected
closed fragment permits neither Mix nor open hypotheses.
\end{hbdefinition}

\begin{hbdefinition}{def:p07-mll-desequentialization}{Construction from a derivation}
The selected rules and auxiliary extensions are displayed on
\cref{sheet:p07-mll-sequents}. A hypothesis leaf produces a single loose
arc. An axiom produces its two-output vertex. A par inference adds a
par vertex to the two specified conclusions of one structure. A tensor
inference takes disjoint copies of its two premise structures and joins
the two specified conclusions to a new tensor vertex. Exchange changes
only the displayed order. Binary Mix is disjoint union; nullary Mix
produces the empty structure. Write \(D(\pi)\) for this construction.
Hypothesis leaves are assumptions of the auxiliary proof, not logical
axioms asserting every formula.

If a proof of \(\Gamma,A\) replaces a marked hypothesis \(A\) of another
open proof, the graph operation identifies that output arc with the
hypothesis arc, retaining the first arc's source and the second's target.
The carried formulas \(\Gamma\) remain separate conclusions.
\end{hbdefinition}

\begin{hbtheorem}{thm:p07-mll-desequentialization}
 {proved=annex:p07.mll.desequentialization}{Derivations produce the stated switching graphs}
Every auxiliary cut-free derivation produces a switching-acyclic structure.
A closed Mix-free derivation produces connected-correct switching graphs.
Graph substitution agrees with proof substitution. For any derivation,
the number of switching components is independent of the switching and is
\[
 d(D(\pi))=1+\#\mathsf{Mix}_2(\pi)-\#\mathsf{Mix}_0(\pi).
\]
\end{hbtheorem}

\section{Girard's original long trips}
\label{sec:p07-mll-original-trips}

\begin{hbdefinition}{def:p07-mll-girard-trip}{Two gates per occurrence}
Girard's cut-free multiplicative proof structure is nonempty, has exactly
one producing link for every occurrence and at most one consuming link,
and uses general-formula axiom, tensor and par links. There are no open
hypotheses in this presentation. Give each occurrence \(F\) the two
exit states \(F^+,F^-\), corresponding to the source's upward and
downward gates. These signs name ports; neither is linear negation.

A setting chooses \(L\) or \(R\) at \emph{every tensor and every par},
independently and for the entire trip. The exact successor clauses
are on \cref{sheet:p07-mll-trips}. They form a permutation of all
\(2p\) exit states, where \(p\) is the number of occurrences. A trip is
an orbit. A long trip contains all \(2p\) states; an orbit of smaller
size is a short trip. Girard's criterion requires every setting to have
one long trip. It is not the modern convention that switches only par
incidences. The distinction includes the source's connectedness
requirement \cite[Definitions~2.1, 2.4--2.5, pp.~29--33]{Girard1987AuthorHosted}.
\end{hbdefinition}

\begin{hbtheorem}{thm:p07-mll-trip-translation}
 {proved=annex:p07.mll.trip-translation}{Original multiplicative derivations have only long trips}
The nonempty structure obtained from a cut-free Mix-free multiplicative
sequent proof by Girard's link construction has one long trip at every
tensor/par setting. General-formula axiom links are allowed.
\end{hbtheorem}

\begin{hbdefinition}{def:p07-mll-girard-empire}{The original tensor-relative empire}
For premises \(A,B\) of a specified tensor link, Girard defines \(eA\)
by a condition over \emph{all} trip settings: both gates of an occurrence
\(C\) must lie in the cyclic interval from \(A^+\) to \(A^-\).
Define \(eB\) symmetrically. This is an intersection of trip intervals,
not the connected component of one arbitrarily chosen switching.
The original argument establishes the passage-order lemmas, seven
closure facts, a setting realizing \(eA\) as one interval, and then
maximal-empire splitting. Not every terminal tensor is a splitting
tensor \cite[pp.~36--40, Lemmas~2.9.1--2.9.4 and Theorems~2.9.5--2.9.7]
{Girard1987AuthorHosted}.
\end{hbdefinition}

\begin{hbtheorem}{thm:p07-mll-original-sequentialization}
 {cited=Girard1987AuthorHosted}{The original trip-based sequentialization theorem}
Every original cut-free multiplicative long-trip net is the link
construction of a Mix-free sequent proof. The tensor case requires
Girard's splitting theorem after terminal pars have been removed
\cite[Theorem~2.9, pp.~35--40]{Girard1987AuthorHosted}.
\end{hbtheorem}

\begin{proofboundary}{def:p07-mll-original-proof-standing}{Original and reconstructed proof routes}
The original trip/sequentialization theorem above retains its cited
standing. The full modern splitting proof below is not a renamed proof
of all of Girard's empire closure cases. In particular the original
seven-case empire argument and its trip-realization/splitting completion
have their distinct full reconstruction in
\cref{sec:p07-girard-trips-empires,thm:p07-girard-empire-sequentialization}.
Nor is the binary reconstruction a
complete reconstruction of the whole original 1989 article's generalized
multiplicative presentation \cite{DanosRegnier1989}. Its exact native
unit, Mix, connectivity and generalized-connective conventions must be
determined from that actual body, not inferred from a later author paper.
\end{proofboundary}

\section{A proved splitting dependency and sequentialization}
\label{sec:p07-mll-sequentialization}

The following proof uses the explicitly different internal-par splitting
method of \cite[Section~3]{DiGuardiaLaurentTortoraVaux2023Sequentialization}.
It is useful because the difficult splitting step is exposed as a
finite graph argument instead of hidden inside a reconstruction algorithm.

\begin{hbdefinition}{def:p07-mll-strong-path}{Bridge, strong path, and splitting par}
A bridge in a path is a consecutive traversal of the two premises of
one par vertex. Its pier is that par vertex. A path is bridge-free when
no such consecutive pair occurs. A nonempty path is strong if, when its
source is a par vertex, its first arc is not a premise of that vertex.
It is strong-weak when it is strong and ends at a par along a premise.
A par is splitting when no simple cycle contains its conclusion arc.

For a vertex \(v\), let \(M_v\) be the set of simple cycles based at
\(v\), containing a conclusion of \(v\), with the least possible number
of bridges among such cycles. It is empty exactly when there is no such
cycle. Define \(v\prec w\) for par vertices when \(w\) is the first
bridge pier on a member of \(M_v\) starting with \(v\)'s conclusion.
These are finite combinatorial definitions, not assumed existence of
a sequentialization.
\end{hbdefinition}

\begin{hblemma}{thm:p07-mll-paths}
 {proved=annex:p07.mll.paths}{The path facts used by splitting}
A strong bridge-free simple cycle is a switching cycle. Concatenating
strong bridge-free paths introduces no bridge at their junction.
A nonempty closed concatenation of simple strong-weak paths contains
a simple cycle as a contiguous subpath.
\end{hblemma}

\begin{hblemma}{thm:p07-mll-bungee}
 {proved=annex:p07.mll.bungee}{The minimal-bridge cycle argument}
Suppose \(\omega\in M_v\), and a simple strong bridge-free path starts
at a bridge pier of \(\omega\) and reaches a vertex of \(\omega\).
Then the structure contains a switching cycle.
\end{hblemma}

\begin{hblemma}{thm:p07-mll-splitting}
 {proved=annex:p07.mll.splitting}{A splitting par exists}
In a switching-acyclic structure, the transitive closure of \(\prec\)
is irreflexive. If the structure contains a par, it contains a splitting
par. This includes disconnected and open structures of the auxiliary
calculus; it does not assume the tensor splitting theorem.
\end{hblemma}

\begin{hbtheorem}{thm:p07-mll-sequentialization}
 {proved=annex:p07.mll.sequentialization}{Full auxiliary sequentialization}
Every switching-acyclic finite structure of
\cref{def:p07-mll-partial-graph} is isomorphic, preserving boundary
occurrences, formula labels and ordered connective ports, to \(D(\pi)\)
for an auxiliary cut-free proof with precisely its marked hypotheses.
\end{hbtheorem}

\begin{hbcorollary}{thm:p07-mll-connected-sequentialization}
 {proved=annex:p07.mll.connected-sequentialization}{The closed connected criterion}
A nonempty closed structure is connected-correct if and only if it is
the construction of a closed cut-free Mix-free proof on the selected
unit-free sheet. The proof removes the auxiliary Mix rules; it does
not silently identify acyclicity with connectedness.
\end{hbcorollary}

\section{Networks that separate the conditions}
\label{sec:p07-mll-net-examples}

\begin{hbexample}{def:p07-mll-valid-invalid}{Two axiom links and two different global outcomes}
Take distinct axiom vertices joining \(a\) to \(a^\perp\), and \(b\)
to \(b^\perp\). Join \(a,b\) to a tensor vertex and
\(a^\perp,b^\perp\) to a par vertex. The conclusions are
\(a\otimes b\) and \(a^\perp\bindnasrepma b^\perp\).
There are two switching graphs. Selecting the \(a^\perp\) premise of
the par leaves the link-node path
\(\mathsf{par},\mathsf{ax}_a,\mathsf{ten},\mathsf{ax}_b\);
selecting the \(b^\perp\) premise reverses the two axiom roles.
All other arcs terminate at boundary points, so both graphs are trees.
A proof is obtained by tensoring the two axiom proofs and then applying
par to their retained dual conclusions.

Replace the par by a tensor, so the second conclusion is now
\(a^\perp\otimes b^\perp\). The four link nodes form an unswitched
cycle through both axioms and both tensors. No choice can remove it.
Thus the identical axiom matching can give a valid or an invalid
network depending on the actual connective links. These are the
occurrence-level variants of Girard's Figure~2
\cite[p.~29]{Girard1987AuthorHosted}.
\end{hbexample}

\begin{hbexample}{def:p07-mll-mix-boundary}{Disconnectedness is a separate failure}
The disjoint union of the two axiom links has no switching cycle.
It has two components and is the image of binary Mix, not a Mix-free
proof. Conversely, joining the two outputs \(a,a^\perp\) of just one
axiom to a tensor produces two parallel unswitched incidences between
that axiom and tensor. This is a two-arc cycle, not a tensor proof.
Planarity and a balanced inventory of atomic labels decide neither case.
\end{hbexample}

\section{Cut rewrites, preservation, and normal forms}
\label{sec:p07-mll-net-cut}

\begin{hbdefinition}{def:p07-mll-net-cut}{Cut links and their two contractions}
In this section a cut structure is closed: every arc has a producing
axiom or connective vertex. A cut vertex has two ordered premise arcs labeled \(A,A^\perp\) and
no conclusion. Its incidences are unswitched, and the cut is symmetric
up to exchanging its premises. Every occurrence still has one producer
and at most one consumer. The switching definition is unchanged.

The principal contraction deletes a cut between a tensor and its dual
par, deletes those two connective vertices and their outgoing arcs,
and installs the two cuts on the corresponding premise pairs.
The axiom contraction deletes the cut and an axiom producing one of
its premises, and joins the axiom's other outgoing arc to the producer
of the opposite cut premise, preserving the former consumer or boundary.
The literal schemas are on \cref{sheet:p07-mll-cut}.
These are incidence surgeries, not substitution of a printed formula
for all equal labels. A cut tying together the two outputs of the
same axiom is excluded by correctness: its two incidences form a cycle.
\end{hbdefinition}

\begin{hbtheorem}{thm:p07-mll-cut-preservation}
 {proved=annex:p07.mll.cut-preservation}{Every contraction preserves correctness}
A contraction of a switching-acyclic cut structure is switching-acyclic
and preserves its switching component count. In particular a
connected-correct cut structure remains connected-correct. No units,
boxes, exponentials or generalized connective contractions are included.
\end{hbtheorem}

\begin{hbtheorem}{thm:p07-mll-cut-termination}
 {proved=annex:p07.mll.cut-termination}{Termination and existence of a cut-free reduct}
Every contraction strictly decreases the number of link vertices.
Every cut in a correct closed structure admits an axiom or principal
contraction. Hence every maximal reduction is finite and ends cut-free.
A nonempty closed correct structure with \(n\) link vertices takes fewer
than \(n\) contractions; the empty structure takes zero.
\end{hbtheorem}

\begin{hbtheorem}{thm:p07-mll-cut-confluence}
 {proved=annex:p07.mll.cut-confluence}{Order independence at the declared scope}
The two contractions are confluent on correct finite closed structures, up to
isomorphism preserving the retained conclusion occurrences, formula
labels and connective port orders. They therefore have a unique
cut-free normal form at that equivalence. This is not uniqueness of
sequent proof trees, and it does not assert a box/exponential normal-form
theorem. Girard's original multiplicative results are
\cite[Propositions~2.10--2.11 and Theorem~2.12, pp.~41--43]
{Girard1987AuthorHosted}.
\end{hbtheorem}

\begin{hbexample}{def:p07-mll-cut-example}{A complete principal-cut computation}
Make two disjoint copies \(N_1,N_2\) of the correct identity network
in \cref{def:p07-mll-valid-invalid}, with roots \(T_i=a\otimes b\)
and \(P_i=a^\perp\bindnasrepma b^\perp\).
Cut \(T_1\) against \(P_2\). Each switching joins the two original
trees through that cut, so it is correct. Its principal contraction
deletes \(T_1,P_2\) and replaces their cut by the two cuts
\((a_1,a_2^\perp)\), \((b_1,b_2^\perp)\).
An axiom contraction on each pair leaves an axiom between
\(a_2,a_1^\perp\), and an axiom between \(b_2,b_1^\perp\).
The remaining roots are \(T_2,P_1\), with exactly the valid identity
network's incidence. The order of the two atomic contractions is immaterial.

Cutting \(T_1\) against \(P_1\) inside just \emph{one} copy is a
different, invalid operation: it adds a path between two vertices already
connected in every switching tree and hence creates a cycle.
The cut rule uses disjoint premise proofs, not arbitrary pairs of
dual conclusions in one proof.
\end{hbexample}

\section{Original scope that the binary theorem does not erase}
\label{sec:p07-mll-source-boundaries}

\begin{proofboundary}{def:p07-mll-box-unit-boundary}{PN0, PN1, PN2, units, and boxes}
Girard's Section~2 begins with the unit-free multiplicative fragment.
Its later notation PN0 includes cut; \(\mathrm{PN0}'\) is cut-free.
PN1 adds unary links and proof boxes; PN2 adds second-order
propositional quantification. In PN1 the unit axiom for \(1\), the
top axiom with auxiliary doors, bottom and quest weakening,
with-box formation, promotion, and quest contraction have their
specific displayed box/door regimes. His generic proof-box switch
is a cyclic permutation of all its nonzero number of doors.
The \(1,\bot\) treatment cannot be obtained by adding a bare nullary
vertex to our no-unit binary graph. Neither can promotion/duplication
be obtained from our size-decreasing cut proof
\cite[pp.~43--47]{Girard1987AuthorHosted}.

The original PN1/PN2 link surfaces and their full box/quantifier
contraction proofs remain separate reconstruction duties. The required
1989 article's full native presentation also remains separate.
Geometry of interaction, including operator execution and its
invariance claims, is the separate N-x-011 source obligation and is not
a theorem of this graph calculus.
\end{proofboundary}

\showcalculus{p07-mll-proof-nets}
\index{proof net!multiplicative}
\index{sequentialization!switching criterion}
\index{cut elimination!proof nets}
